%========================================
% MatinBook - Stage 11: Cross-References Test
% Test: Hyperref + Cleveref + All Reference Commands
% Purpose: Validate intelligent cross-referencing system
%========================================

\makeatletter
\def\input@path{{../}}
\makeatother

\documentclass{matinbook}

\title{آزمون سیستم ارجاع هوشمند}
\author{تیم توسعه MatinBook}
\date{تابستان ۱۴۰۵}

\begin{document}

%----------------------------------------
% Title Page
%----------------------------------------
\maketitle

%----------------------------------------
% Chapter 1: Creating Reference Targets
%----------------------------------------
\chapter{ایجاد اهداف ارجاع}
\label{chap:targets}

\section{معادلات}

چند معادله معروف که به آن‌ها ارجاع خواهیم داد:

\begin{equation}
    E = mc^{2}
    \label{eq:einstein}
\end{equation}

\begin{equation}
    a^{2} + b^{2} = c^{2}
    \label{eq:pythagoras}
\end{equation}

\begin{equation}
    e^{i\pi} + 1 = 0
    \label{eq:euler}
\end{equation}

\section{قضایا و تعاریف}

\begin{theorem}[قضیه مقدار میانی]
    اگر تابع $f$ روی بازه بسته $[a, b]$ پیوسته باشد و
    $y$ عددی بین $f(a)$ و $f(b)$ باشد، آنگاه
    حداقل یک $c \in [a, b]$ وجود دارد که $f(c) = y$.
    \label{thm:intermediate-value}
\end{theorem}

\begin{definition}[پیوستگی]
    تابع $f$ در نقطه $x_{0}$ \textbf{پیوسته} است اگر:
    \[
        \lim_{x \to x_{0}} f(x) = f(x_{0})
    \]
    \label{def:continuity}
\end{definition}

\begin{lemma}[لم فشردگی]
    هر دنباله کراندار در $\R$ دارای زیردنباله همگراست.
    \label{lem:bolzano-weierstrass}
\end{lemma}

\begin{corollary}
    هر دنباله یکنوا و کراندار همگراست.
    \label{cor:monotone-convergence}
\end{corollary}

\begin{proposition}
    مجموعه اعداد گویا $\Q$ در $\R$ چگال است.
    \label{prop:q-dense}
\end{proposition}

\section{مثال‌ها و نکات}

\begin{example}[تابع پیوسته]
    تابع $f(x) = x^{2}$ روی تمام $\R$ پیوسته است،
    زیرا برای هر $x_{0}$:
    \[
        \lim_{x \to x_{0}} x^{2} = x_{0}^{2} = f(x_{0})
    \]
    \label{ex:continuous-function}
\end{example}

\begin{remark}
    عکس قضیه \thmref{thm:intermediate-value} لزوماً برقرار نیست.
    یعنی ممکن است تابعی در شرط مقدار میانی صدق کند
    اما پیوسته نباشد.
    \label{rem:ivt-converse}
\end{remark}

\section{شکل و جدول}

\begin{figure}[h]
\centering
\begin{latin}
\begin{tikzpicture}
  \begin{axis}[
    width=10cm, height=5cm,
    xlabel={$x$},
    ylabel={$f(x)$},
    title={Continuous Function},
    grid=major,
    domain=-2:2,
    samples=100
  ]
    \addplot[matinblue, thick] {x^2};
    \addplot[matinred, only marks, mark=*] coordinates {(1, 1)};
    \node[above] at (axis cs:1,1) {$(1, f(1))$};
  \end{axis}
\end{tikzpicture}
\end{latin}
\caption{نمودار تابع $f(x) = x^{2}$}
\label{fig:parabola}
\end{figure}

\begin{table}[h]
\centering
\caption{مقادیر تابع $f(x) = x^{2}$}
\label{tab:parabola-values}
\begin{tabular}{|c|c|}
\hline
\textbf{$x$} & \textbf{$f(x)$} \\
\hline
$-2$ & $4$ \\
$-1$ & $1$ \\
$0$ & $0$ \\
$1$ & $1$ \\
$2$ & $4$ \\
\hline
\end{tabular}
\end{table}

%----------------------------------------
% Chapter 2: Testing All Reference Commands
%----------------------------------------
\chapter{آزمون تمام دستورات ارجاع}
\label{chap:references}

\section{ارجاع استاندارد با \lr{\textbackslash ref}}

\begin{itemize}
    \item معادله اینشتین: \ref{eq:einstein}
    \item قضیه مقدار میانی: \ref{thm:intermediate-value}
    \item نمودار سهمی: \ref{fig:parabola}
\end{itemize}

\section{ارجاع هوشمند با \lr{\textbackslash cref}}

دستور \texttt{\textbackslash cref} به صورت خودکار
نوع عنصر را تشخیص می‌دهد و پیشوند مناسب اضافه می‌کند:

\begin{itemize}
    \item \cref{eq:einstein} معادله معروف اینشتین است
    \item طبق \cref{thm:intermediate-value}، تابع پیوسته
    شرط مقدار میانی را دارد
    \item در \cref{fig:parabola} نمودار سهمی را می‌بینیم
    \item \cref{tab:parabola-values} مقادیر تابع را نشان می‌دهد
    \item \cref{def:continuity} تعریف پیوستگی است
    \item \cref{ex:continuous-function} یک مثال است
    \item \cref{rem:ivt-converse} یک نکته مهم است
    \item \cref{lem:bolzano-weierstrass} لم فشردگی است
    \item \cref{cor:monotone-convergence} نتیجه لم است
    \item \cref{prop:q-dense} گزاره چگالی $\Q$ است
\end{itemize}

\section{ارجاع اختصاصی \lr{MatinBook}}

این دستورات برای راحتی بیشتر تعریف شده‌اند:

\begin{itemize}
    \item \textbf{معادله:} \meqref{eq:euler} زیباترین معادله ریاضی است
    \item \textbf{قضیه:} \thmref{thm:intermediate-value} پایه آنالیز است
    \item \textbf{تعریف:} \defref{def:continuity} مفهوم حد را دقیق می‌کند
    \item \textbf{لم:} \lemref{lem:bolzano-weierstrass} در اثبات‌ها کاربرد دارد
    \item \textbf{نتیجه:} \corref{cor:monotone-convergence} از لم نتیجه می‌شود
    \item \textbf{مثال:} \exref{ex:continuous-function} ساده ولی گویاست
    \item \textbf{شکل:} \figref{fig:parabola} یک سهمی ساده را نشان می‌دهد
    \item \textbf{جدول:} \tabref{tab:parabola-values} مقادیر عددی را فهرست کرده
    \item \textbf{فصل:} در \chref{chap:targets} اهداف ارجاع را تعریف کردیم
    \item \textbf{بخش:} \secref{sec:smart-refs} ارجاع‌های هوشمند را توضیح می‌دهد
\end{itemize}

\section{ارجاع‌های هوشمند در عمل}
\label{sec:smart-refs}

حالا بیاییم از ترکیب این ارجاع‌ها در یک متن واقعی استفاده کنیم:

\medskip

طبق \thmref{thm:intermediate-value}، اگر تابعی پیوسته باشد
(\defref{def:continuity})، آنگاه شرط مقدار میانی را دارد.
\exref{ex:continuous-function} نشان می‌دهد که $x^{2}$
(\figref{fig:parabola}) این شرط را دارد.
اما \ref{rem:ivt-converse} هشدار می‌دهد که عکس
\thmref{thm:intermediate-value} برقرار نیست.
برای مطالعه بیشتر، به \chref{chap:cross-chapter} مراجعه کنید.

%----------------------------------------
% Chapter 3: Range References
%----------------------------------------
\chapter{ارجاع‌های بازه‌ای}
\label{chap:range-refs}

\section{ارجاع به چند معادله}

قابلیت \lr{cleveref} امکان ارجاع به چند عنصر را
به صورت فشرده فراهم می‌کند:

\begin{itemize}
    \item \cref{eq:einstein,eq:pythagoras,eq:euler}
    سه معادله معروف هستند
    \item طبق قضایای \cref{thm:intermediate-value,lem:bolzano-weierstrass,cor:monotone-convergence}
    نتایج مهمی در آنالیز داریم
\end{itemize}

\section{ارجاع بازه‌ای}

با استفاده از \texttt{\textbackslash crefrange}:

\begin{itemize}
    \item معادلات \crefrange{eq:einstein}{eq:euler}
    همگی از معادلات بنیادی هستند
\end{itemize}

%----------------------------------------
% Chapter 4: Cross-Chapter References
%----------------------------------------
\chapter{ارجاع بین فصلی}
\label{chap:cross-chapter}

\section{ارجاع به فصل‌های دیگر}

حالا می‌خواهیم به عناصری که در \chref{chap:targets}
تعریف کردیم، از این فصل ارجاع دهیم:

\begin{itemize}
    \item معادله \meqref{eq:pythagoras} که در
    \chref{chap:targets} تعریف شد، اساس هندسه است
    \item \thmref{thm:intermediate-value} (از \chref{chap:targets})
    یکی از قضایای اساسی آنالیز ریاضی است
    \item \figref{fig:parabola} در \chref{chap:targets}
    نمودار تابع درجه دو را نشان می‌دهد
\end{itemize}

\section{ارجاع به این فصل از فصول قبلی}

همچنین می‌توان از فصول قبلی به این فصل ارجاع داد.
برای مثال، اگر در \chref{chap:targets} بودیم،
می‌توانستیم بنویسیم:
«در \chref{chap:cross-chapter} ارجاع بین فصلی
را بررسی خواهیم کرد.»

%----------------------------------------
% Summary
%----------------------------------------
\chapter{خلاصه نتایج}

\section{وضعیت سیستم ارجاع}

\begin{table}[h]
\centering
\caption{وضعیت تمام دستورات ارجاع}
\label{tab:reference-status}
\begin{tabular}{|c|C{4cm}|c|}
\hline
\textbf{دستور} & \textbf{کاربرد} & \textbf{وضعیت} \\
\hline
\lr{\textbackslash ref} & شماره عنصر بدون پیشوند & \textcolor{matingreen}{✅} \\
\lr{\textbackslash cref} & شماره با پیشوند هوشمند & \textcolor{matingreen}{✅} \\
\lr{\textbackslash crefrange} & ارجاع بازه‌ای & \textcolor{matingreen}{✅} \\
\lr{\textbackslash meqref} & ارجاع به معادله & \textcolor{matingreen}{✅} \\
\lr{\textbackslash thmref} & ارجاع به قضیه & \textcolor{matingreen}{✅} \\
\lr{\textbackslash defref} & ارجاع به تعریف & \textcolor{matingreen}{✅} \\
\lr{\textbackslash lemref} & ارجاع به لم & \textcolor{matingreen}{✅} \\
\lr{\textbackslash corref} & ارجاع به نتیجه & \textcolor{matingreen}{✅} \\
\lr{\textbackslash exref} & ارجاع به مثال & \textcolor{matingreen}{✅} \\
\lr{\textbackslash figref} & ارجاع به شکل & \textcolor{matingreen}{✅} \\
\lr{\textbackslash tabref} & ارجاع به جدول & \textcolor{matingreen}{✅} \\
\lr{\textbackslash chref} & ارجاع به فصل & \textcolor{matingreen}{✅} \\
\lr{\textbackslash secref} & ارجاع به بخش & \textcolor{matingreen}{✅} \\
\hline
\end{tabular}
\end{table}

\section{نتیجه}

\textbf{سیستم ارجاع هوشمند \lr{MatinBook} نسخه ۱ با موفقیت آزمایش شد.}
تمامی ۱۳ دستور ارجاع شامل \lr{hyperref}، \lr{cleveref}
و دستورات اختصاصی فارسی به درستی کار می‌کنند. 🎉

\end{document}